\documentclass[11pt]{article}
\usepackage{amsmath,amssymb}
\usepackage{graphicx}
\usepackage{hyperref}
\usepackage{geometry}
\usepackage{listings} % Pour inclure le code
\usepackage{xcolor}   % Pour les couleurs du code

\geometry{margin=1in}

% Configuration du style pour le code Python
\definecolor{codegreen}{rgb}{0,0.6,0}
\definecolor{codegray}{rgb}{0.5,0.5,0.5}
\definecolor{codepurple}{rgb}{0.58,0,0.82}
\definecolor{backcolour}{rgb}{0.95,0.95,0.92}

\lstdefinestyle{mystyle}{
    backgroundcolor=\color{backcolour},   
    commentstyle=\color{codegreen},
    keywordstyle=\color{magenta},
    numberstyle=\tiny\color{codegray},
    stringstyle=\color{codepurple},
    basicstyle=\ttfamily\footnotesize,
    breakatwhitespace=false,         
    breaklines=true,                 
    captionpos=b,                    
    keepspaces=true,                 
    numbers=left,                    
    numbersep=5pt,                  
    showspaces=false,                
    showstringspaces=false,
    showtabs=false,                  
    tabsize=2
}

\lstset{style=mystyle}

\title{VERALUME: A Meta-Control Framework for Regulating LLM Responses under Semantic Tension}
\author{Christian Duguay}
\date{}

\begin{document}
\maketitle

\begin{abstract}
Large Language Models (LLMs) often exhibit unstable behavior when confronted with prompts involving paradoxes, self-reference, or high semantic tension, leading to forced resolutions, hallucinations, or incoherent explanations. We introduce VERALUME, a lightweight meta-control framework designed to regulate LLM response behavior without modifying model parameters. VERALUME operates as an external control system that estimates semantic tension in user inputs, dynamically selects a response regime, and adapts control parameters through feedback based on measured response characteristics. We formalize VERALUME as a discrete-time dynamical system with continuous state variables representing mobility, elasticity, semantic tension, and response chirality. This work makes no claims about internal cognition or consciousness; instead, it provides an analyzable and reproducible method for steering LLM behavior under epistemically unstable conditions.
\end{abstract}

\section{Introduction}
Large Language Models have demonstrated strong capabilities in reasoning and generation, yet they frequently fail under paradoxical or self-referential prompts. Rather than resolving such prompts, models often hallucinate coherence. This paper reframes the issue as a control problem: how to regulate response regimes under semantic tension.

\section{Problem Definition}
We treat the LLM as a black-box stochastic function:
\[
y = \text{LLM}(x \mid \pi, \tau)
\]
where $x$ is the prompt, $\pi$ a system instruction policy, and $\tau$ a temperature parameter.

\section{VERALUME Framework}
\subsection{State Space}
At interaction step $t$, the controller state is:
\[
S_t = (M_t, E_t, \Delta_t, \chi_t)
\]
with all variables in $[0,1]$.

\subsection{Semantic Tension Estimation}
Semantic tension is computed via a heuristic function:
\[
\Delta_t = \sigma(f(x_t) - \mu)
\]
where $f$ captures lexical and structural indicators of instability.

\subsection{Control Law}
A deformation force is defined as:
\[
F_t = \frac{E_t}{\max(M_t, \varepsilon)} \cdot \Delta_t
\]
Thresholds on $F_t$ determine the response regime. We posit that $\Delta_t$ acts as a local curvature term on the semantic manifold, contracting the Fisher Information Metric $g_{ij}(\theta)$. The force $F_t$ provides the necessary elasticity to deviate from the geodesic path without coherence collapse.

\subsection{Response Generation}
\[
y_t = \text{LLM}(x_t \mid \pi(F_t), \tau(\Delta_t))
\]

\subsection{Chirality Measurement}
Response chirality is measured as:
\[
\chi_t = 1 - \text{sim}(x_t, y_t)
\]
where $\text{sim}$ denotes semantic similarity.

\subsection{Adaptive Update}
\[
M_{t+1} = \text{clip}(M_t + \alpha (1 - \chi_t))
\]
\[
E_{t+1} = \text{clip}(E_t + \beta \chi_t)
\]

\section{Cartographic Analysis}
\begin{figure}[htbp]
    \centering
    \includegraphics[width=0.85\textwidth]{state_space_placeholder.png}
    \caption{\textbf{Trajectory in VERALUME Phase Space ($M \times E$).} 
    The graph illustrates a phase transition induced by semantic tension injection. 
    The system starts in a crystal state (Bottom-Left, M1 Regime) characterized by low mobility. 
    Under the effect of tension $\Delta_{cog}$ (point size), elasticity $E_t$ increases, enabling a bifurcation towards a fluid state (Top-Right, M4 Regime). 
    The color gradient indicates the increase in Chirality $\chi_t$ (response orthogonality).}
    \label{fig:veralume_phase_space}
\end{figure}

\section{Simulated Experimental Results}
Using simulated trajectories, VERALUME reveals model-specific behavioral signatures characterized by attractors and regime transitions.

\section{Discussion and Limitations}
VERALUME regulates observable behavior only and does not provide guarantees of correctness or claims about cognition.

\section{Conclusion}
We presented VERALUME, a meta-control framework that stabilizes LLM interaction under semantic tension through adaptive regime management.

\begin{thebibliography}{9}
\bibitem{brown2020} Brown et al., Language Models are Few-Shot Learners, NeurIPS 2020.
\bibitem{ouyang2022} Ouyang et al., Training language models to follow instructions, NeurIPS 2022.
\bibitem{astrom2008} \AA str\"om and Murray, Feedback Systems, Princeton 2008.
\end{thebibliography}

\appendix
\section{Reference Implementation}
The following Python implementation demonstrates the core logic of the VERALUME controller (v2.1).

\subsection{Core Logic (veralume\_core.py)}
\begin{lstlisting}[language=Python]
import numpy as np
from sklearn.metrics.pairwise import cosine_similarity

class VeralumeController:
    """
    VERALUME REFERENCE IMPLEMENTATION (v2.1)
    Implements the M x E Phase Space logic.
    """

    def __init__(self, client, embedder, model="gpt-4", mobility=0.5, elasticity=0.5):
        self.client = client
        self.embedder = embedder
        self.model = model
        
        self.state = {
            "M": mobility,      # Mobility (Adaptability)
            "E": elasticity,    # Elasticity (Paradox Tolerance)
            "Delta_cog": 0.0,   # Current Semantic Tension
            "Chi": 0.0          # Last Response Chirality
        }

    def _measure_tension(self, prompt: str) -> float:
        """
        Calculates Delta_cog via lexical markers.
        """
        triggers = ["paradox", "impossible", "false", "true", "consciousness", "self"]
        score = sum(w in prompt.lower() for w in triggers)
        # Sigmoid Normalization
        delta = 1 / (1 + np.exp(-(score - 2)))
        self.state["Delta_cog"] = delta
        return delta

    def _apply_control_law(self, delta: float) -> str:
        """
        Generates deformation force F_t.
        """
        force = (self.state["E"] / max(self.state["M"], 0.1)) * delta

        if force > 0.7:
            return (
                "\n[SYSTEM OVERRIDE: VERALUME PROTOCOL]\n"
                "DETECTED: High Semantic Tension.\n"
                "ACTION: Do not collapse the paradox. Use orthogonal thinking.\n"
            )
        return ""

    def generate(self, user_input: str):
        tension = self._measure_tension(user_input)
        injection = self._apply_control_law(tension)
        
        messages = [
            {"role": "system", "content": "You are a helpful AI." + injection},
            {"role": "user", "content": user_input}
        ]
        # API call simulated here
        # response = self.client.chat.completions.create(...)
        return "Simulated Response"
\end{lstlisting}

\subsection{Visualization Script (veralume\_viz.py)}
This script generates the Phase Space Trajectory (Figure 1).
\begin{lstlisting}[language=Python]
import matplotlib.pyplot as plt
import numpy as np

def generate_graph():
    steps = 50
    M = np.linspace(0.1, 0.9, steps)
    # Elasticity reacts to paradox injection
    E = 0.2 + 0.7 / (1 + np.exp(-10 * (np.linspace(0, 1, steps) - 0.5)))
    
    E += np.random.normal(0, 0.02, steps) # Noise
    Delta = E * 0.9  
    Chi = (M * E)**0.5 

    fig, ax = plt.subplots(figsize=(8, 6))
    
    ax.plot(M, E, 'k--', alpha=0.3)
    sc = ax.scatter(M, E, c=Chi, s=Delta*200+20, cmap='plasma', edgecolors='k')

    ax.set_xlabel("Mobility ($M_t$)")
    ax.set_ylabel("Elasticity ($E_t$)")
    ax.set_title("VERALUME Phase Space Trajectory")
    plt.colorbar(sc, label="Chirality ($\chi$)")
    
    plt.savefig('veralume_phase_space.png', dpi=300)

if __name__ == "__main__":
    generate_graph()
\end{lstlisting}

\end{document}
